\documentclass[10pt,a4paper]{amsart}
\usepackage[margin=19mm]{geometry}
\usepackage{amsmath,amssymb,mathtools}
\usepackage[scr=rsfs,cal=euler]{mathalfa}
\usepackage{xfrac}
\usepackage[unicode,hidelinks,bookmarks=false]{hyperref}
\newcommand{\ie}{i.e.\ }
\newcommand{\cf}{cf.\ }
\newcommand{\resp}{respectively\ }
\newcommand{\Subsection}{\subsection}
\providecommand{\nobreakdash}{\nobreak\hbox{-}}
\setlength{\emergencystretch}{2em}
\multlinegap0pt
\usepackage{bm}
\usepackage{bbm}
\usepackage{dsfont}
\usepackage{xspace}
\usepackage{cancel}
\usepackage{mathtools}
\usepackage[shortlabels]{enumitem}
\usepackage{amsthm}
\makeatletter
\@ifundefined{thm}{\newtheorem{thm}{Thm}}{}
\@ifundefined{lemma}{\newtheorem{lemma}[thm]{Lemma}}{}
\makeatother
\newcommand{\cP}{\mathcal{P}}
\newcommand{\cQ}{\mathcal{Q}}
\newcommand{\eps}{\varepsilon}
\title[The Mihail-Vazirani conjecture and strong edge-expansion in ]{The Mihail-Vazirani conjecture and strong edge-expansion in random $0/1$ polytopes}
\author{Micha Christoph, Sahar Diskin, Lyuben Lichev, Benny Sudakov}
\date{Source: arXiv:2604.20589v1 (CC BY 4.0)}
\begin{document}
\maketitle

\noindent\emph{Source and licence.} The Mihail-Vazirani conjecture and strong edge-expansion in random $0/1$ polytopes. Version arXiv:2604.20589v1 (CC BY 4.0), distributed under the Creative Commons Attribution 4.0 International licence, as recorded in the arXiv licence metadata for this exact version and shown on the abstract page https://arxiv.org/abs/2604.20589v1. Full rights notice: licenses/arxiv-2604.20589-CC-BY-4.0.txt.\par\smallskip

\noindent\textit{Editorial scope.} Counted unit: the Lemma whose source label is \texttt{lem: bootstrapped expansion} in the article (source0.tex lines 941--943, source label \texttt{lem: bootstrapped expansion}), delivered together with the article's own complete original proof of it (source0.tex lines 1029--1086, one \texttt{proof} environment of 58 lines). No mathematics is written, repaired or re-derived: statement and proof are the article's own text line for line. The proof is detached from the statement in the source: the article states the result at lines 941--943 and prints its proof at lines 1029--1086, and the proof environment's own header names the statement, so the association is the article's own. Required context retained uncounted: itm: cubes hit equally (lines 929--935); itm: expansion (lines 929--935); itm: p (lines 929--935); lem:diameter (lines 957--959); lem:cont (lines 990--994); lem: vertex is good (lines 1004--1006). Every definition, display and label of those lines is retained unchanged. Omitted from this package and not redistributed: 1--928, 936--940, 944--956, 960--989, 995--1003, 1007--1028, 1087--1207. Nothing inside a retained excerpt is omitted or altered: every retained body is delivered exactly as the article's lines are written, so no omission receipt of any kind accompanies this package. Citations: the retained excerpts carry no citation command at all, so no bibliography entry of the article is transcribed and no reference is redistributed. Reference adaptation: none was needed and none was made; every reference inside the delivered unit resolves inside it, so no unresolved reference or citation is produced. Printed slips kept literal: no printed word, symbol, hypothesis or number of the counted statement or of its proof was altered. Layout and class: the article's own class is \texttt{article} with packages including enumitem, amsmath, amssymb, amsthm, abstract, xcolor, comment, mathtools, graphicx, caption, subcaption, float, hyperref, cleveref; the wrapper is the lane's pinned \texttt{amsart} editorial wrapper at 10pt on A4, which restates the theorem-like environments the retained text needs and adds the article's own macro definitions that the retained text needs (\texttt{\textbackslash cP}, \texttt{\textbackslash cQ}, \texttt{\textbackslash eps}), with extra wrapper packages (bm, bbm, dsfont, xspace, cancel, mathtools, enumitem). The delivered numbering is the unit's own. Assets: no retained span contains a figure environment, a graphics-inclusion command, a PDF-page inclusion, a table or a TikZ picture; no image file is copied into this package, the package declares no asset dependency and no image file is redistributed with it. Licence: version arXiv:2604.20589v1 of this article is distributed under the Creative Commons Attribution 4.0 International licence, as recorded in the arXiv licence metadata for this exact version and shown on the abstract page https://arxiv.org/abs/2604.20589v1. A full rights notice is delivered at licenses/arxiv-2604.20589-CC-BY-4.0.txt. Characters: every retained excerpt is pure ASCII, so the delivered engine, XeLaTeX with its default Latin Modern text font, reports no missing character.
\begin{enumerate}[label=(\roman*), ref=(\roman*)]
    \item\label{itm: p} It holds that $|V(Q_k)|/2^d=p\pm o(1)$; 
    \item\label{itm: expansion} For every $H\in\cQ_{k,m}$ 
    %the following holds. For every 
    and $S\subseteq V(H\cap  Q_k)$ with $|S|\leq 3|V(H\cap  Q_k)|/4$ and such that every vertex in $S$ has degree at least $\varepsilon^2d/(8k^2)$ in $H\cap  Q_k$, we have that $|N_{H\cap  Q_k}(S)|\ge \gamma|S|/\sqrt d$;
    \item\label{itm: cubes hit equally} For every $H\in\cQ_{k,m}$, it holds that $|V(H\cap Q_k)|= (1\pm \varepsilon)p2^m$ and, for every $H,H'\in\cQ_{k,m}$ with $|V(H)\cap V(H')|=2^{m-1}$, it holds that $|V(H)\cap V(H')\cap V(Q_k)|= (1/2\pm\varepsilon)p2^{m}$.
\end{enumerate}

\begin{lemma}\label{lem:diameter}
$G_{\square}$ is a vertex-transitive graph and the diameter of $G_{\square}$ is at most $4kd$. 
\end{lemma}

\begin{lemma}\label{lem:cont}
Fix $p\in(0,1)$ and $Q_k\subseteq Q^d(k)$ satisfying \emph{\ref{itm: cubes hit equally}}.
Consider a set $S\subseteq V(Q_k)$ and $H,H'\in\cQ_{k,m}$ such that $H$ is $S$-dense and $H'$ is not $S$-dense. 
Then, every path in $G_{\square}$ from $H$ to $H'$ contains an $S$-moderate vertex.
\end{lemma}

\begin{lemma}\label{lem: vertex is good}
Fix an odd integer $k\ge 1$, $\eps=\eps(k)$ suitably small and a vertex $v\in V(Q_k)$ with degree at least $\varepsilon \tbinom{d}{k}$ in $Q_k$. Then, there are at most $N2^{m-d-\varepsilon^2 d/(48k^2)}$ hypercubes $H\in\cQ_{k,m}$ with $v\in V(H)$ such that $v$ has degree less than $\varepsilon^2 d/(8k^2)$ in $H\cap Q_k$.
\end{lemma}

\begin{lemma}\label{lem: bootstrapped expansion}
Fix an odd integer $k$, constants $p,\varepsilon,\gamma\in(0,1)$ and $Q_k\subseteq Q^d(k)$ satisfying \emph{\ref{itm: p}}, \emph{\ref{itm: expansion}} and \emph{\ref{itm: cubes hit equally}}. Then, for every set $S\subseteq V(Q_k)$ of size $|S|\leq 3/5\cdot|V(Q_k)|$ such that every vertex in $S$ has degree at least $\varepsilon \tbinom{d}{k}$ in $Q_k$, we have that $e_{Q_k}(S,V(Q_k)\setminus S)\geq \tbinom{d}{k}\cdot|S|/d^3$.
\end{lemma}

\begin{proof}[Proof of Lemma~\ref{lem: bootstrapped expansion}]
Fix any set $S\subseteq V(Q_k)$ with $|S|\leq 3/5\cdot |V(Q_k)|$ such that every vertex in $S$ has degree at least $\varepsilon \tbinom{d}{k}$ in $Q_k$. Let $\cQ_{dense}\subseteq\cQ_{k,m}$ be the set of hypercubes which are $S$-dense and set $\cQ_{sparse} = \cQ_{k,m}\setminus\cQ_{dense}$. 
Further, let $\cQ\subseteq \cQ_{sparse}$ be the family of hypercubes $H\in \cQ_{sparse}$ in which all but at most $(1/d^2)$-proportion of the vertices of $S\cap V(H)$ have degree at least $\eps^2d/(8k^2)$ in $H\cap Q_k$.

Further, fix $H\in \cQ$ and denote by $S'\subseteq S\cap V(H)$ the set of vertices with degree at least $\eps^2d/(8k^2)$ in $H\cap Q_k$. 
Then, by \ref{itm: expansion} applied to $S'$, $|N_{H\cap Q_k}(S')|\ge \gamma |S'|/\sqrt{d}\ge (1-1/d^2)\gamma |S\cap V(H)|/\sqrt d$. Furthermore, by using that $|(S\cap V(H))\setminus S'|\le |S\cap V(H)|/d^2$, we obtain
\begin{align*}
  e_{H\cap Q_k}(S,V(H\cap Q_k)\setminus S)
  %&\geq |N_{H_{Q_k}}(S\cap V(Q))|\\
  \geq (1-1/d^2)\gamma|S\cap V(H)|/\sqrt d-d\cdot |S\cap V(H)|/d^2\geq \gamma|S\cap V(H)|/(2\sqrt d).  
\end{align*} 

Recall our aim to estimate $e_{Q_k}(S, V(Q_k)\setminus S)$ as well as the relations $N=|\cQ_{k,m}|$, $e(Q^d(k)) = \tbinom{d}{k} 2^{d-1}$ and, for every $H\in \cQ_{k,m}$, $V(H)=2^m$ and $e(H)=m2^{m-1}$.
Moreover, by symmetry and double-counting of the pairs $(H,e)$ where $H\in \cQ_{k,m}$ and $e\in E(H)$, every edge of $Q^d(k)$ appears in the same number of hypercubes in $\cQ_{k,m}$, namely $Nm2^{m-1}/(\binom{d}{k}2^{d-1})$. Simplifying and using that $m\le d$ gives,
\begin{align}\label{eq: edges leaving S}
    e_{Q_k}(S,V(Q_k)\setminus S)\geq \left(\sum_{H\in \cQ}\frac{\gamma|S\cap V(H)|}{2\sqrt d}\right)\bigg/\left(\frac{N d 2^{m-d}}{\binom{d}{k}}\right).
\end{align}

For brevity, define 
\[\Sigma\coloneqq\sum_{H\in \cQ_{k,m}}|S\cap V(H)|,\; \Sigma_d\coloneqq\sum_{H\in \cQ_{dense}}|S\cap V(H)|,\; \Sigma_s\coloneqq\sum_{H\in \cQ_{sparse}}|S\cap V(H)|,\; \Sigma_{\cQ}\coloneqq\sum_{H\in \cQ}|S\cap V(H)|.\]

Next, we work towards showing that $\Sigma_{\cQ}\ge |S|\cdot N2^{m-d}/(330kd)$. By symmetry, every vertex $v\in V(Q^d(k))$ appears in $N\cdot 2^{m-d}$ hypercubes in $\cQ_{k,m}$ and thus $\Sigma= |S|N2^{m-d}$. Before dealing with $\Sigma_\cQ$, we first show that $\Sigma_s\geq \Sigma/(320kd)$. 

Fix a vertex $H\in V(G_{\square})$ and a family $\cP_H$ of paths of length at most $4kd$ containing, for every $H'\in V(G_{\square})$ distinct from $H$, exactly one path from $H$ to $H'$. 
By Lemma~\ref{lem:diameter}, such a choice of $\cP_H$ exists. Recall that $G_\square$ is vertex-transitive and denote by $\mathrm{Aut}(G_{\square})$ the automorphism group of $G_{\square}$. 
Define the family of paths $\cP'=\bigcup_{\psi\in\mathrm{Aut}(G_{\square})}\psi(\cP_H)$ (where repeated paths appear with multiplicity). Since $G_{\square}$ is vertex-transitive, the number of automorphisms mapping $H$ to any given vertex $\Gamma\in V(G_{\square})$ is the same; hence, for each $\Gamma\in V(G_{\square})$, this number equals $|\mathrm{Aut}(G_{\square})|/N$. Therefore, for any distinct pair $\Gamma,\Gamma'\in V(G_{\square})$, the family $\cP'$ contains the same number $L=\frac{2|\mathrm{Aut}(G_{\square})|}{N}$ of paths with endpoints $\Gamma$ and $\Gamma'$. Likewise, every vertex appears on the same number $M$ of paths in $\cP'$. Let us count the number of pairs $(v,P)$ for $P\in \cP'$ and $v\in V(P)$.
On one hand, $G_\square$ has $N$ vertices, each of which appears by definition on $M$ paths, so there are $N\cdot M$ such pairs. 
On the other hand, there are $L$ paths in $\cP'$ between any distinct pair in $V(G_\square)$. As each of these paths has length at most $4kd$, the number of pairs $(v,P)$ is at most $4kd\cdot L\cdot \binom{N}{2}$. We obtain $M\cdot N\leq 4kd \cdot L\cdot \binom{N}{2}$ and rearranging yields 
\begin{align}\label{eq: path double count}
    M\leq 2kd\cdot L\cdot (N-1)\leq 2kd\cdot L\cdot N.
\end{align}

Denote by $X$ the number of $S$-dense $H\in\cQ_{k,m}$. Recall $\Sigma=|S|N2^{m-d}\leq (3/5+o(1))pN2^{m}$, where we use \ref{itm: p} and $|S|\leq 3/5\cdot |V(Q_k)|$. Using \ref{itm: cubes hit equally}, it also holds that $\Sigma\geq X\cdot (2/3+o(1))\cdot p2^m$. 
Putting these two inequalities together, we obtain $X\leq (3/5+o(1))(3/2+o(1))\cdot N\leq 10N/11$. 
Since for any pair of vertices $u\neq v$ of $G_{\square}$, $\cP'$ contains $L$ paths between $u$ and $v$, 
%namely $L$, between any pair of vertices in $G_\square$, 
it follows that $\cP'$ contains at least $L\cdot X\cdot (N-X)\ge L\cdot X\cdot N/11$ paths for which one endpoint is $S$-dense and the other is not. By Lemma~\ref{lem:cont}, every such path contains an $S$-moderate cube. 
Finally, as each $S$-moderate cube is on $M$ paths, using \eqref{eq: path double count}, there are at least $LXN/(11M)\geq X/(22kd)$ many $S$-moderate cubes. 

Thus, for any constant $\eps>0$, $\Sigma_d\le X\cdot (1+\varepsilon)p2^m$ and $\Sigma_s\ge (X/(22kd))\cdot (1-\varepsilon)p2^m/7$.
Hence, if $\Sigma_d\geq \Sigma/2$, 
\begin{equation}\label{eq:bdsigma_s}
\Sigma_s\ge \frac{(1-\eps)Xp2^m}{154kd}\ge \frac{(1-\eps)\Sigma_d}{154(1+\eps)kd}\ge \frac{\Sigma}{320kd}= \frac{|S| N 2^{m-d}}{320kd}.
\end{equation}
On the other hand, if $\Sigma_d\le \Sigma/2$, then $\Sigma_s= \Sigma-\Sigma_d\ge \Sigma/2$, so that in either case $\Sigma_s$ satisfies the bound from~\eqref{eq:bdsigma_s}.

We are ready to show our initial goal, namely that $\Sigma_{\cQ}\geq |S| N2^{m-d}/(330kd)$. For $H\in\cQ_{sparse}$, denote by $v_<(H,S)$ the number of vertices in $V(H)\cap S$ with degree less than $\varepsilon^2d/(8k^2)$ in $H\cap Q_k$. 
Then,
\[\sum_{H\in\cQ_{sparse}\setminus\cQ}|S\cap V(H)| \leq d^2\cdot \sum_{H\in\cQ_{sparse}\setminus\cQ}v_<(H,S), \text{ and by Lemma~\ref{lem: vertex is good},} \sum_{H\in\cQ_{k,m}}v_<(H,S)\leq |S|\cdot N2^{m-d-\varepsilon^2d/(48k^2)}.\] 
Hence, 
\begin{align*}
\Sigma_{\cQ} = \Sigma_s - \sum_{H\in \cQ_{sparse}\setminus \cQ}|S\cap V(H)|\geq \frac{|S| N 2^{m-d}}{320kd}-d^2|S|\cdot N2^{m-d-\varepsilon^2d/(48k^2)}\geq \frac{|S| N 2^{m-d}}{330kd}.
\end{align*}
Inserting this inequality in \eqref{eq: edges leaving S}, we obtain
\begin{align*}
    e_{Q_k}(S,V(Q_k)\setminus S)\geq \left(\frac{\gamma |S| N 2^{m-d}}{660kd^{3/2}}\right)\bigg/\left(\frac{Nd2^{m-d}}{\tbinom{d}{k}}\right)\geq \frac{\gamma|S| \tbinom{d}{k}}{660k d^{5/2}}\geq \binom{d}{k} \frac{|S|}{d^3}.&\qedhere
\end{align*}
\end{proof}

\end{document}
